\chapter{Практика}


\section{Настройка интерфейсов роутера}

\begin{lstlisting}[language=C, caption={Пример настройки интерфейсов роутера}]
#R1
configure

interface gigabitethernet 1/0/1
ip address 192.168.100.1/30
ip firewall disable
exit

interface gigabitethernet 1/0/2
ip address 172.16.100.1/30
ip firewall disable
exit

1interface gigabitethernet 1/0/3
ip address 192.168.0.1/23
ip firewall disable
exit

do commit
do confirm

#R2
configure

interface gigabitethernet 1/0/1
ip address 10.10.100.2/30
ip firewall disable
exit

interface gigabitethernet 1/0/2
ip address 192.168.100.2/30
ip firewall disable
exit

1interface gigabitethernet 1/0/3
ip address 10.10.0.1/23
ip firewall disable
exit

do commit
do confirm

#R3
configure

interface gigabitethernet 1/0/1
ip address 172.16.100.2/30
ip firewall disable
exit

interface gigabitethernet 1/0/2
ip address 10.10.100.1/30
ip firewall disable
exit

interface gigabitethernet 1/0/3
ip address 172.16.0.1/23
ip firewall disable
exit

do commit
do confirm
\end{lstlisting}

\section{Настройка ip на ПК}

\begin{lstlisting}[language=C, caption={Пример настройки ip на ПК}]
#VPC1
ip 192.168.1.0/23 192.168.0.1
#VPC2
ip 10.10.1.0/23 10.10.0.1
#VPC3
ip 172.16.1.0/23 172.16.0.1
\end{lstlisting}

\section{Создание процесса OSPF}

\begin{lstlisting}[language=C, caption={Пример создания процесса OSPF}]
# создаем ospf процесс с id=250
Router-left(config)# router ospf 250
# задаем id маршрутизатора в этом процессе
Router-left(config-ospf)# router-id 1.1.1.1
# задаем ospf область
Router-left(config-ospf)# area 0.0.0.0

# применяем настройки
Router-left(config-ospf-area)# enable
Router-left(config-ospf-area)# exit
Router-left(config-ospf)# enable

# На R2,R3 настройка аналогична, но router-id изменится на 2.2.2.2 и 3.3.3.3 соответсвенно
\end{lstlisting}

\section{Добавление интерфейсов в процесс ospf}

\begin{lstlisting}[language=C, caption={Пример добавления интерфейсов в процесс ospf}]
Router-left(config)# interface gigabitethernet 1/0/1-3
# привязка интерфейсов к ospf процессу
Router-left(config-if-gi)# ip ospf instance 250
# привязка интерфейсов к ospf области
Router-left(config-if-gi)# ip ospf area 0.0.0.0
# включение ospf на интерфейсах
Router-left(config-if-gi)# ip ospf
Router-left(config-if-gi)# exit

# На R2,R3 настройки аналогичны
\end{lstlisting}

\section{Перевод интерфейсов в пассивный режим}

На каждом роутере интерфейс gi1/0/3 подключен к оконечному устройству, который не участвует в OSPF процессе, но в ходе работы ospf он будет получать пакеты,
которые никак использовать не будет. Получаем лишний трафик + это не безопасно. Проблема решается путем перевода интерфейса в пассивный режим.

\begin{lstlisting}[language=C, caption={Пример перевода интерфейсов в пассивный режим}]
Router-left(config)# interface gigabitethernet 1/0/3
Router-left(config-if-gi)# ip ospf passive-interface

# На R2,R3 настройки аналогичны
\end{lstlisting}


\section{Проверка таблицы маршрутизации}

\begin{figure}[H] 
    \centering
    \includegraphics[width=0.75\textwidth]{rt.png}
    \caption{Таблица маршрутизации R1}
\end{figure}

Видим новые маршруты с типом O, полученные через OSPF. 

\section{Проверка сетевой связности}

Сделаем trace ПК1-ПК3

\begin{figure}[H] 
    \centering
    \includegraphics[width=0.75\textwidth]{trace.png}
    \caption{Трассировка маршрута}
\end{figure}

Видно, что пакеты идут по маршруту R1-R3.\\

Теперь будем выполнять пинг ПК1-ПК3 и удалим канал между R1-R3

\begin{figure}[H] 
    \centering
    \includegraphics[width=0.75\textwidth]{ping.png}
    \caption{Пинг с удалением канала}
\end{figure}

В отличие от RIP, OSPF не нужно время на сходимость сети после изменения топологии, что позволяет не терять пакеты.
Можно заметить, что в какой-то момент получение эхо-ответа от пункта назначения становится больше, именно в этот момент и произошло
изменение пути. В конце сделал trace, который подтверждает, что маршрут изменился и теперь трафик идет по маршруту R1-R2-R3










\endinput